% ICML 2026 final-project report.
% Drop-in: replace \usepackage{icml2026} with the actual ICML 2026 style file
% provided by the conference. Tested with the latest icml2025.sty as a placeholder.
\documentclass{article}

% ICML 2026 style placeholder. The conference provides icml2026.sty;
% drop it next to this file and uncomment the line below. Until then, basic
% article + booktabs + hyperref keeps the document compilable.
% \usepackage[accepted]{icml2026}

\usepackage{microtype}
\usepackage{graphicx}
\usepackage{booktabs}
\usepackage{amsmath, amssymb}
\usepackage{hyperref}
\usepackage{url}
\usepackage{xcolor}
\usepackage{enumitem}
\usepackage{caption}
\usepackage{subcaption}
\usepackage[margin=1in]{geometry}

\title{Pixel Reconstruction vs.\ Latent Prediction for Self-Supervised
Representation Learning of Active Matter Dynamics}
\author{Person A \quad Person B \\ NYU CSCI-GA 2572 (Spring 2026)}
\date{}

\begin{document}
\maketitle

\begin{abstract}
We study self-supervised representation learning for the
\texttt{active\_matter} subset of \emph{The Well} \citep{ohana2025well} under
the strict frozen-encoder evaluation protocol mandated by our course (single
linear probe and $k$-Nearest Neighbours regression for the physical parameters
$\alpha$ and $\zeta$). We compare two SSL paradigms on a
\emph{byte-identical} 3D ConvNeXt encoder so that any difference in downstream
mean-squared error (MSE) is causally attributable to the SSL objective alone:
(i) \textbf{SIGReg-JEPA} \citep{balestriero2026lejepa} -- latent prediction with
distributional regularization toward an isotropic Gaussian; and (ii) a
\textbf{VideoMAE/SimMIM-hybrid} \citep{tong2022videomae,xie2022simmim} masked
autoencoder that we adapt for ConvNeXt by replacing masked tubes with a
learnable mask token. We further train an end-to-end supervised baseline on the
same encoder as an upper-bound anchor, and report a representation-collapse,
isotropy and rank analysis suite that applies the LeJEPA Epps--Pulley statistic
as a \emph{diagnostic} for any encoder. Our results
\textcolor{red}{[TODO: fill in headline numbers when training+eval complete]}.
\end{abstract}

\section{Introduction}
\label{sec:intro}

Self-supervised learning has emerged as a foundation paradigm for representation
learning in vision, language and, more recently, the physical sciences
\citep{qu2026spatiotemporal,balestriero2023cookbook}. In this work we focus on
\emph{active matter}, a subset of \emph{The Well} \citep{ohana2025well} that
simulates rod-like active particles in a viscous Stokes fluid via continuum
kinetic theory. Each trajectory is governed by two latent physical parameters
-- the active dipole strength $\alpha$ and the steric alignment $\zeta$ -- and
the central question is whether a self-supervised encoder, trained
\emph{without access to those parameters}, learns a representation from which
$\alpha$ and $\zeta$ can be recovered by a simple frozen-encoder probe.

The closest published work to ours is \citet{qu2026spatiotemporal}, who report
that on this exact dataset a JEPA-style encoder beats a VideoMAE-style encoder
by ${\sim}51\%$ relative MSE (0.16 $\to$ 0.08) when the encoder is followed by
a 100-epoch \emph{attentive-probe} fine-tuning stage. The course rules
explicitly prohibit attentive probes and any fine-tuning of the backbone; only
a single linear layer and $k$-NN are allowed for evaluation. Our
\textbf{contribution} is to ask the same JEPA-vs-MAE question under this much
stricter protocol, while \emph{controlling for the encoder family} (both
methods share the byte-identical 3D ConvNeXt encoder, isolating the SSL
objective as the only causal variable).

In short, our contributions are:
\begin{itemize}[topsep=2pt,itemsep=2pt]
    \item A \textbf{controlled comparison} of two SSL objectives -- SIGReg
    latent prediction (Person~A's stream) vs.\ VideoMAE/SimMIM pixel
    reconstruction (Person~B's stream) -- on \emph{the same} 3D ConvNeXt
    encoder, evaluated under the strict frozen-linear+kNN protocol.
    \item A \textbf{SimMIM adaptation for non-ViT encoders}: depthwise 3D
    ConvNeXts cannot use VideoMAE's asymmetric encoder-decoder, so we replace
    masked tubes with a learnable per-channel mask token (Sec.\ \ref{sec:method-pixel}).
    \item A \textbf{representation-quality analysis suite} that applies the
    LeJEPA Epps--Pulley statistic \citep{balestriero2026lejepa} as a
    \emph{diagnostic} for any frozen encoder, alongside effective rank,
    participation ratio, condition number and channel/frame ablations
    (Sec.\ \ref{sec:analysis}).
    \item An \textbf{end-to-end supervised baseline} trained on the same
    encoder, providing an upper-bound anchor and an interesting
    "frozen-supervised" diagnostic.
\end{itemize}

\section{Setup}
\label{sec:setup}

\subsection{Dataset}
The \texttt{active\_matter} subset comprises 8\,750 / 1\,200 / 1\,300 sliding
windows of 16 frames at $224{\times}224$ spatial resolution and 11 physical
channels (concentration scalar, 2-channel velocity, 4-channel orientation
tensor, 4-channel strain-rate tensor); 45 unique $(\alpha,\zeta)$ combinations
total. Detailed statistics are in Tab.~\ref{tab:dataset} (Appendix).

\subsection{Encoder (shared between both SSL streams)}
We use a 3D ConvNeXt encoder with stages of dimensions
$(32, 64, 128, 256, 256)$ and per-stage block depths $(2, 2, 4, 8, 2)$, a
$2{\times}4{\times}4$ stem patch and 0.05 stochastic depth. The total spatial
downsampling is $32\times$ and temporal downsampling is $16\times$, yielding a
$(B, 256, 7, 7)$ bottleneck per clip. The encoder has 7.7\,M parameters.

\subsection{Evaluation protocol}
For every encoder we extract frozen features and run \emph{both}
(i) a single $\mathbf{nn.Linear}(256, 2)$ probe trained with MSE loss against
$z$-scored $(\alpha, \zeta)$ on the train split, and (ii) a $k$-Nearest
Neighbours regressor with hyperparameters swept over
$k\in\{1,3,5,10,20,50\}$, weights $\in\{$uniform, distance$\}$, metrics
$\in\{$cosine, euclidean$\}$ and feature normalization
$\in\{$z-score, $\ell_2$, z+$\ell_2\}$. Model selection uses the validation
split; the test split is touched only once per encoder for the final number.
Label normalization (z-score) is fit on \emph{train labels only}.

\section{SSL Objectives}
\label{sec:method}

\subsection{Latent prediction: SIGReg-JEPA (Person A)}
\label{sec:method-jepa}
\textcolor{red}{[Person A: fill in the SIGReg / LeJEPA section -- the
Epps--Pulley distributional regularizer formulated against $\mathcal{N}(0, I)$,
context/target temporal split, predictor architecture, etc.]}

\subsection{Pixel reconstruction: VideoMAE/SimMIM hybrid (Person B)}
\label{sec:method-pixel}
We adapt VideoMAE \citep{tong2022videomae} for our 3D ConvNeXt encoder. Three
adaptations are needed compared to the canonical ViT-based recipe:
\begin{enumerate}[topsep=2pt,itemsep=2pt]
    \item \textbf{Mask token replacement (SimMIM-style).} Depthwise 3D
    convolutions require the full input tensor; we cannot use VideoMAE's
    asymmetric encoder-decoder which drops masked tokens. Following SimMIM
    \citep{xie2022simmim} we instead replace each masked tube with a learnable
    per-channel mask token, broadcast over the tube extent, and feed the full
    tensor to the encoder.
    \item \textbf{Tube size aligned with encoder downsampling.} Per SimMIM
    Sec.~4.1.2, the optimal mask patch size matches the encoder stride. With
    $32\times$ spatial and $16\times$ temporal downsampling, our tube size is
    $(16, 32, 32)$, giving a $7{\times}7=49$ tube grid per sample.
    \item \textbf{Lightweight prediction head.} SimMIM Tab.~2 shows lightweight
    heads transfer better than heavy decoders. We use a 5-stage
    transposed-conv decoder (350\,k parameters) that mirrors the encoder
    downsampling, mapping $(256, 7, 7) \to (11, 16, 224, 224)$. Total model
    size is 8.07\,M parameters, well under the assignment's 100\,M cap.
\end{enumerate}

The reconstruction target is per-tube z-scored pixel values (VideoMAE
Sec.~3.3), and the loss is MSE over masked positions only (SimMIM Tab.~4
shows prediction-only beats full reconstruction):
\begin{equation}
\mathcal{L}_{\text{MAE}} = \frac{1}{|\mathcal{M}|}
    \sum_{(c,t,h,w) \in \mathcal{M}}
    \left( \hat{x}_{c,t,h,w} - \tilde{x}_{c,t,h,w} \right)^2,
\end{equation}
where $\mathcal{M}$ is the set of masked pixel coordinates and $\tilde{x}$ is
the per-tube z-scored target. The encoder is trained for 25 epochs with AdamW
($\text{lr}=1.5{\times}10^{-4}$, weight decay 0.05), a cosine schedule with
2-epoch linear warmup, and BF16 autocast.

\subsection{Supervised baseline (Person B)}
\label{sec:method-supervised}
We also train an end-to-end supervised baseline on the \emph{same} 3D ConvNeXt
encoder followed by an $\mathbf{nn.Linear}(256, 2)$ head, optimizing $z$-scored
MSE directly against $(\alpha, \zeta)$. This serves as an upper-bound anchor
and exposes how much information about the physical parameters is recoverable
from the encoder family at all. Trained for 30 epochs with AdamW
($\text{lr}=3{\times}10^{-4}$).

\section{Results}
\label{sec:results}

\textcolor{red}{[TODO: regenerate Tab.~\ref{tab:main-results} from
videomae/artifacts/figures/figures\_summary.json once eval completes.]}

\begin{table}[t]
\centering
\caption{Active matter parameter prediction. Z-scored MSE (lower is better) on
the test split with a frozen encoder, evaluated with both a single linear probe
and $k$-NN regression. Best per row in \textbf{bold}.}
\label{tab:main-results}
\small
\begin{tabular}{lrrr}
\toprule
\textbf{Encoder} & \textbf{Params (M)} & \textbf{Linear probe} & \textbf{kNN} \\
\midrule
Supervised (ours)             & 7.7 & TBD & TBD \\
VideoMAE / SimMIM mask=0.60   & 7.7 & TBD & TBD \\
VideoMAE / SimMIM mask=0.75   & 7.7 & TBD & TBD \\
SIGReg-JEPA (Person A)        & 7.7 & TBD & TBD \\
\midrule
\multicolumn{4}{l}{\emph{External references (different protocol -- attentive probe FT for 100 epochs):}} \\
JEPA, ConvNeXt-tiny           & ${\sim}3.3$ & 0.079 & --- \\
VideoMAE, ViT-tiny/16         & ${\sim}5.5$ & 0.160 & --- \\
\bottomrule
\end{tabular}
\end{table}

\section{Analysis}
\label{sec:analysis}

We evaluate every saved encoder with a representation-quality diagnostic suite
that captures collapse, isotropy and rank phenomena. All metrics operate on the
frozen-encoder embedding dump produced by \texttt{export\_embeddings.py}.

\paragraph{Collapse via per-dim std.} The sorted per-dimension standard
deviation curve (Fig.~\ref{fig:collapse}) makes representation collapse visible:
a collapsed encoder exhibits a heavy long tail with most dimensions near zero,
while a healthy encoder shows roughly uniform std.

\paragraph{Effective rank and participation ratio.} We compute
$\exp(H(p))$ where $p$ are the normalized eigenvalues of the embedding
covariance, plus the participation ratio
$(\sum \lambda_i)^2 / \sum \lambda_i^2$. Both are widely used proxies for the
true dimensionality of a representation manifold.

\paragraph{Condition number and off-diagonal correlation.} The covariance
condition number $\lambda_{\max} / \lambda_{\min}$ and the mean absolute
off-diagonal correlation reveal feature redundancy and ill-conditioning that
hurts both linear probes and kNN with Euclidean distances.

\paragraph{Epps--Pulley distance to $\mathcal{N}(0, I)$.} We repurpose the
LeJEPA SIGReg statistic \citep{balestriero2026lejepa} as a generic diagnostic:
the average squared characteristic-function error along random projections,
weighted by $e^{-t^2/2}$ and integrated to $t=3$. This is precisely Person A's
training objective (for SIGReg-JEPA), but applied here to \emph{any} encoder's
embeddings as a goodness-of-fit-to-isotropic-Gaussian metric.

\paragraph{Channel and frame ablations.} \textcolor{red}{[TODO: add results once
ablation runs complete.]}

\section{Limitations and Ethics}

We trained on a single seed per configuration due to compute budget; multiple
seeds would tighten the error bars on our headline comparison. We did not
train a from-scratch ViT-based VideoMAE that would precisely replicate Qu et
al.'s baseline -- our VideoMAE uses the colleague's ConvNeXt encoder for
causal cleanliness of the SSL-objective comparison, at the cost of being
slightly different from the published recipe. The dataset itself is a synthetic
fluid simulation; no human or animal data are involved.

\paragraph{Compute accounting.} All experiments ran on 2x NVIDIA B200
180\,GB GPUs, BF16 autocast, on a single rented node. End-to-end:
\textcolor{red}{[TODO: total GPU-hours and peak VRAM]}.

\section*{Statement of Contributions}

\textbf{Person A} authored the \texttt{physrl/} package (3D ConvNeXt encoder,
JEPA latent-prediction objective, VICReg and SIGReg loss functions, the
training/eval/sweep CLIs reused unchanged here, and the related-work review of
JEPA-style SSL). Person~A's stream covers Sec.~\ref{sec:method-jepa} and the
SIGReg-JEPA rows of Tab.~\ref{tab:main-results}.

\textbf{Person B} authored the \texttt{videomae/} package (the
VideoMAE/SimMIM-hybrid trainer, end-to-end supervised baseline,
representation-quality analysis suite, the figure regeneration script, the
detached training launchers, the dataset hash and ENV.md/requirements lock for
reproducibility) and the corresponding paper sections
(\ref{sec:method-pixel},~\ref{sec:method-supervised},~\ref{sec:analysis}).
Files in \texttt{videomae/} whose names also appear in \texttt{physrl/} are
derivative copies with rewritten imports; modifications are documented in their
module docstrings.

% References
\begin{thebibliography}{9}

\bibitem[Qu et al.(2026)]{qu2026spatiotemporal}
Helen Qu, Rudy Morel, Michael McCabe, Alberto Bietti, Fran\c{c}ois Lanusse,
Shirley Ho, Yann LeCun.
\newblock Representation Learning for Spatiotemporal Physical Systems.
\newblock \emph{ICLR 2026 Workshop on AI \& PDE}, arXiv:2603.13227, 2026.

\bibitem[Tong et al.(2022)]{tong2022videomae}
Zhan Tong, Yibing Song, Jue Wang, Limin Wang.
\newblock VideoMAE: Masked Autoencoders are Data-Efficient Learners for
Self-Supervised Video Pre-Training.
\newblock \emph{NeurIPS 2022}, arXiv:2203.12602.

\bibitem[Xie et al.(2022)]{xie2022simmim}
Zhenda Xie, Zheng Zhang, Yue Cao, Yutong Lin, Jianmin Bao, Zhuliang Yao,
Qi Dai, Han Hu.
\newblock SimMIM: A Simple Framework for Masked Image Modeling.
\newblock \emph{CVPR 2022}, arXiv:2111.09886.

\bibitem[Balestriero \& LeCun(2026)]{balestriero2026lejepa}
Randall Balestriero, Yann LeCun.
\newblock LeJEPA: Provable and Scalable Self-Supervised Learning Without the
Heuristics.
\newblock arXiv:2511.08544, 2026.

\bibitem[Bardes et al.(2021)]{bardes2021vicreg}
Adrien Bardes, Jean Ponce, Yann LeCun.
\newblock VICReg: Variance-Invariance-Covariance Regularization for
Self-Supervised Learning.
\newblock arXiv:2105.04906, 2021.

\bibitem[Liu et al.(2022)]{liu2022convnext}
Zhuang Liu, Hanzi Mao, Chao-Yuan Wu, Christoph Feichtenhofer, Trevor Darrell,
Saining Xie.
\newblock A ConvNet for the 2020s.
\newblock \emph{CVPR 2022}.

\bibitem[Ohana et al.(2025)]{ohana2025well}
Ruben Ohana et al.
\newblock The Well: a Large-Scale Collection of Diverse Physics Simulations
for Machine Learning.
\newblock arXiv:2412.00568, 2025.

\bibitem[Balestriero et al.(2023)]{balestriero2023cookbook}
Randall Balestriero et al.
\newblock A Cookbook of Self-Supervised Learning.
\newblock arXiv:2304.12210, 2023.

\bibitem[He et al.(2022)]{he2022mae}
Kaiming He, Xinlei Chen, Saining Xie, Yanghao Li, Piotr Doll\'ar, Ross Girshick.
\newblock Masked Autoencoders Are Scalable Vision Learners.
\newblock \emph{CVPR 2022}, arXiv:2111.06377.

\end{thebibliography}

\end{document}
